\begin{thebibliography}{99}
\bibitem{ref1} Alexander Alder. Grenzrang und Grenzkomplexität aus algebraischer und topologischer Sicht. PhD thesis, Zentralstelle der Studentenschaft, 1984.
\bibitem{ref2} Eric Allender, Peter Bürgisser, Johan Kjeldgaard-Pedersen, and Peter Miltersen. On the Complexity of Numerical Analysis. SIAM Journal on Computing, 38, January 2006. Preliminary version in the 21st Annual IEEE Conference on Computational Complexity (CCC’06). doi: 10.1109/CCC.2006.30.
\bibitem{ref3} Robert Andrews. Algebraic Hardness Versus Randomness in Low Characteristic. In 35th Computational Complexity Conference (CCC 2020), volume 169 of Leibniz International Proceedings in Informatics (LIPIcs), pages 37:1–37:32. Schloss Dagstuhl–Leibniz-Zentrum für Informatik, 2020. doi:10.4230/LIPIcs.CCC.2020.37.
\bibitem{ref4} Jean Berstel and Christophe Reutenauer. Rational Series and Their Languages. Springer- Verlag, Berlin, Heidelberg, 1988. URL: https://dl.acm.org/doi/book/10.5555/52107.
\bibitem{ref5} Markus Bläser and Gorav Jindal. On the Complexity of Symmetric Polynomials. In 10th Innovations in Theoretical Computer Science Conference (ITCS’19), volume 124 of LIPIcs, pages 47:1–47:14. Schloss Dagstuhl - Leibniz-Zentrum für Informatik, 2019. doi:10.4230/ LIPIcs.ITCS.2019.47.
\bibitem{ref6} Alin Bostan, Gilles Christol, and Philippe Dumas. Fast computation of the Nth term of an algebraic series over a finite prime field. In Proceedings of the ACM on International Symposium on Symbolic and Algebraic Computation (ISSAC’16), pages 119–126, 2016. doi: 10.1145/2930889.2930904.
\bibitem{ref7} Karl Bringmann, Christian Ikenmeyer, and Jeroen Zuiddam. On Algebraic Branching Programs of Small Width. J. ACM, 65(5):1–29, 2018. (Preliminary version in the 32nd Computational Complexity Conference (CCC’17). doi:10.1145/3209663.
\bibitem{ref8} Peter Bürgisser. The complexity of factors of multivariate polynomials. Foundations of Computational Mathematics, 4(4):369–396, 2004. arXiv:1812.06828.
\bibitem{ref9} Peter Bürgisser. On defining integers and proving arithmetic circuit lower bounds. Computa- tional Complexity, 18(1):81–103, 2009. doi:10.1007/s00037-009-0260-x.
\bibitem{ref10} Peter Bürgisser, Michael Clausen, and Amin Shokrollahi. Algebraic complexity theory, volume 315. Springer Science \& Business Media, 2013.
\bibitem{ref11} David V Chudnovsky and Gregory V Chudnovsky. On expansion of algebraic functions in power and Puiseux series, I. Journal of Complexity, 2(4):271–294, 1986. URL: https: //www.sciencedirect.com/science/article/pii/0885064X86900063.
\bibitem{ref12} Michael Coons. The Transcendence of Series Related to Stern’s Diatomic Sequence. Interna- tional Journal of Number Theory, 06, November 2011. doi:10.1142/S1793042110002958.
\bibitem{ref13} Michael Coons and Peter Borwein. Transcendence of power series for some number theoretic functions. Proceedings of the American Mathematical Society, 137, July 2008. doi:10.1090/ S0002-9939-08-09737-2.
\bibitem{ref14} Wellington De Melo and Benar Fux Svaiter. The cost of computing integers. Proceedings- American Mathematical Society, 124:1377–1378, 1996. URL: https://www.ams.org/journals/ proc/1996-124-05/S0002-9939-96-03173-5/S0002-9939-96-03173-5.pdf.
\bibitem{ref15} Richard A. Demillo and Richard J. Lipton. A probabilistic remark on algebraic program testing. Information Processing Letters, 7(4):193–195, 1978. URL: https://www.sciencedirect.com/ science/article/abs/pii/0020019078900674.
\bibitem{ref16} Pranjal Dutta. Real tau-Conjecture for sum-of-squares: A unified approach to lower bound and derandomization. In 16th International Computer Science Symposium in Russia (CSR 2021), 2021. URL: https://drive.google.com/file/d/1X8eo9GM4SCNsC2vWjPbUwMX0vff5i2k3/ view.
\bibitem{ref17} Pranjal Dutta, Nitin Saxena, and Amit Sinhababu. Discovering the roots: Uniform closure results for algebraic classes under factoring. In Proceedings of the 50th Annual ACM SIGACT Symposium on Theory of Computing, pages 1152–1165, 2018. URL: https://www.cse.iitk. ac.in/users/nitin/papers/factor-closure.pdf.
\bibitem{ref18} Pranjal Dutta, Nitin Saxena, and Thomas Thierauf. A Largish Sum-Of-Squares Implies Circuit Hardness and Derandomization. In 12th Innovations in Theoretical Computer Science Conference (ITCS 2021), volume 185 of Leibniz International Proceedings in Informatics (LIPIcs), pages 23:1–23:21. Schloss Dagstuhl–Leibniz-Zentrum für Informatik, 2021. URL: 10.4230/LIPIcs.ITCS.2021.23.
\bibitem{ref19} Joshua A Grochow et al. Complexity in ideals of polynomials: questions on algebraic complexity of circuits and proofs. Bulletin of EATCS, 2(130), 2020. URL: http://bulletin.eatcs.org/ index.php/beatcs/article/view/607.
\bibitem{ref20} Joshua A. Grochow, Ketan D. Mulmuley, and Youming Qiao. Boundaries of VP and VNP. In 43rd International Colloquium on Automata, Languages, and Programming (ICALP 2016), volume 55, pages 34:1–34:14, 2016. URL: https://core.ac.uk/download/pdf/62922137.pdf.
\bibitem{ref21} Gorav Jindal. On approximate polynomial identity testing and real root finding. PhD thesis, Saarland University, 2019. doi:10.22028/D291-29880.
\bibitem{ref22} Erich Kaltofen. Uniform closure properties of p-computable functions. In Proceedings of the eighteenth annual ACM symposium on Theory of computing, pages 330–337, 1986. doi: 10.1145/12130.12163.
\bibitem{ref23} Erich Kaltofen. Single-factor Hensel lifting and its application to the straight-line complexity of certain polynomials. In Proceedings of the 19th annual ACM symposium on Theory of computing (STOC’87), pages 443–452, 1987. doi:10.1145/28395.28443.
\bibitem{ref24} Pascal Koiran. Valiant’s model and the cost of computing integers. computational complexity, 13(3):131–146, 2005.
\bibitem{ref25} Pascal Koiran. Shallow circuits with high-powered inputs. Innovations in Computer Sci- ence (ICS), 2011. URL: https://hal-ens-lyon.archives-ouvertes.fr/ensl-00477023v4/ document.
\bibitem{ref26} Pascal Koiran and Sylvain Perifel. Interpolation in Valiant’s theory. Computational Complexity, 20(1):1–20, 2011. doi:10.1007/s00037-011-0002-8.
\bibitem{ref27} Pascal Koiran, Natacha Portier, Sébastien Tavenas, and Stéphan Thomassé. A $\tau$ -Conjecture for Newton Polygons. Foundations of computational mathematics, 15(1):185–197, 2015. doi: 10.1007/s10208-014-9216-x.
\bibitem{ref28} FV Kuhlmann. On convergent power series, 1996. URL: https://www.mathi.uni-heidelberg. de/~roquette/KONVPOTREIHEN.pdf.
\bibitem{ref29} Hsiang Kung and Joseph Traub. All Algebraic Functions Can Be Computed Fast. J. ACM,
\bibitem{ref30} Dick Lipton and Ken Regan. Factoring and factorials, February 2009. URL: https://rjlipton. wordpress.com/2009/02/23/factoring-and-factorials/.
\bibitem{ref31} Richard J Lipton. Polynomials with 0-1 coefficients that are hard to evaluate. SIAM Journal on Computing, 7(1):61–69, 1978. Preliminary version in the 16th Annual Symposium on Foundations of Computer Science (FOCS 1975). URL: https://epubs.siam.org/doi/abs/ 10.1137/0207004?journalCode=smjcat.
\bibitem{ref32} Richard J Lipton. Straight-line complexity and integer factorization. In International Algorithmic Number Theory Symposium (ANTS 94), pages 71–79. Springer, 1994. doi: 10.1007/3-540-58691-1\_45.
\bibitem{ref33} Richard J Lipton and Larry J Stockmeyer. Evaluation of polynomials with super- preconditioning. Journal of Computer and System Sciences, 16(2):124–139, 1978. URL: https://www.sciencedirect.com/science/article/pii/0022000078900417.
\bibitem{ref34} Meena Mahajan. Algebraic Complexity Classes. In Perspectives in Computational Complexity, pages 51–75. Springer, 2014. doi:10.1007/978-3-319-05446-9\_4.
\bibitem{ref35} Carlos Moreira. On asymptotic estimates for arithmetic cost functions. Proceedings of the American Mathematical Society, 125(2):347–353, 1997. URL: https://www.jstor.org/ stable/2161660.
\bibitem{ref36} Ketan D Mulmuley. The GCT program toward the P vs. NP problem. Communications of the ACM, 55(6):98–107, 2012. doi:10.1145/2184319.2184341.
\bibitem{ref37} Danny Nguyen and Igor Pak. Complexity of short generating functions. In Forum of Mathematics, Sigma, volume 6. Cambridge University Press, 2018. arXiv:1702.08660.
\bibitem{ref38} Øystein Ore. Über höhere kongruenzen. Norsk Mat. Forenings Skrifter, 1(7):15, 1922.
\bibitem{ref39} Igor Pak. Complexity problems in enumerative combinatorics. In Proceedings of the Inter- national Congress of Mathematicians – Rio de Janeiro 2018. Vol. IV. Invited lectures, pages 3153–3180. World Sci. Publ., Hackensack, NJ, 2018. doi:10.1142/9789813272880\_0176.
\bibitem{ref40} David A Plaisted. New NP-hard and NP-complete polynomial and integer divisibility problems. Theoretical Computer Science, 31(1-2):125–138, 1984. Preliminary in the 17th Annual Symposium on Foundations of Computer Science (FOCS 1976). URL: https: //www.sciencedirect.com/science/article/pii/0304397584901300.
\bibitem{ref41} Jacob T Schwartz. Fast probabilistic algorithms for verification of polynomial identities. Journal of the ACM (JACM), 27(4):701–717, 1980. doi:10.1145/322217.322225.
\bibitem{ref42} Adi Shamir. Factoring numbers in O (logn) arithmetic steps. Information Processing Let- ters, 8(1):28–31, 1979. URL: https://www.sciencedirect.com/science/article/abs/pii/ 0020019079900875.
\bibitem{ref43} Amir Shpilka and Amir Yehudayoff. Arithmetic Circuits: A survey of recent results and open questions. Foundations and Trends® in Theoretical Computer Science, 5(3–4):207–388, 2010. doi:10.1561/0400000039.
\bibitem{ref44} Michael Shub and Steve Smale. On the intractability of Hilbert’s Nullstellensatz and an algebraic version of “NP $\neq$ P?”. Duke Math. J., 81(1):47–54 (1996), 1995. A celebration of John F. Nash, Jr. doi:10.1215/S0012-7094-95-08105-8.
\bibitem{ref45} Volker Strassen. Vermeidung von divisionen. Journal für die reine und angewandte Mathematik, 264:184–202, 1973.
\bibitem{ref46} L Valiant. Reducibility by algebraic projections in: Logic and algorithmic. In Symposium in honour of Ernst Specker, pages 365–380, 1982.
\bibitem{ref47} Leslie G Valiant. Completeness classes in algebra. In Proceedings of the 11th Annual ACM symposium on Theory of computing, pages 249–261. ACM, 1979. doi:10.1145/800135.804419.
\bibitem{ref48} Joachim Von Zur Gathen and Jürgen Gerhard. Modern computer algebra. Cambridge university press, 2013.
\bibitem{ref49} Wact. Some Accessible Open Problems. Workshop on Algebraic Complexity Theory (WACT 2016).   URL: https://www.cs.tau.ac.il/~shpilka/wact2016/concreteOpenProblems/ openprobs.pdf.
\bibitem{ref50} Klaus W Wagner. The complexity of combinatorial problems with succinct input representation. Acta informatica, 23(3):325–356, 1986. doi:10.1007/BF00289117.
\bibitem{ref51} Richard Zippel. Probabilistic Algorithms for Sparse Polynomials. In Proceedings of the International Symposium on Symbolic and Algebraic Computation, EUROSAM ’79, pages 216–226, 1979. doi:10.1007/3-540-09519-5\_73.
\bibitem{ref52} J.von zur Gathen and V. Strassen. Some polynomials that are hard to compute. Theoretical Computer Science, 11(3):331–335, 1980. URL: http://www.sciencedirect.com/science/ article/pii/0304397580900201.
\end{thebibliography}
